\section{Limitations}
\label{sec:limitations}

We state these as findings about what the evidence supports, not as disclaimers.

\paragraph{The benchmark is synthetic, and the fraud in it is invented.} It supports comparisons
between methods under identical conditions. It does not support statements about real fraud rates,
real detection performance or any institution's customers. Every figure here is a property of this
generator.

\paragraph{No fraud parameter is sourced.} Of \nFraudParams{} fraud parameters,
\nFraudParamsChoice{} are modelling choices and \nFraudParamsCalibrated{} are calibrated to
specification targets. No publication read gives fraud incidence by type for these markets. Any
result that depends on the level or mix of fraud, rather than on the relative ordering of methods,
inherits this assumption. The sourced parameters (\nParamsSourced{} of \nParamsTotal{}) fix
population and volume structure, partly from a neighbouring market's statistics.

\paragraph{The benchmark is burst-separable.} Its fraud is injected as bursts, so a single velocity
feature reaches \nFloorM{}, four feature groups each nearly solve it alone, leave-one-country-out
and the novel variant are both null, and the one pre-registered non-burst variant is almost
invisible to the model. Detection of burst-structured fraud is what the benchmark measures; nothing
here speaks to fraud with other structure, except as the failure in
Section~\ref{sec:generalisation}. The specification's distribution targets (overall and test-period
fraud rate, channel and country mix) are requirements met by construction, not observations.

\paragraph{No real-data validation.} None has been done. A validation on real data would need to
show, at minimum: the single-feature floor on real engineered features (whether real fraud is as
burst-separable as this benchmark's); the ensemble's margin over that floor; calibration in the
decision region; leave-one-country-out on real, country-labelled fraud; and detection of
documented non-burst typologies such as the reversal scam. It would require data access agreements,
a data protection impact assessment, and hosting in a permitted location (Section~\ref{sec:ethics}).

\paragraph{Scale and volume.} Results are at about one million rows, not the five million the
specification requires (\notyet{M10}), and the models were trained on \nTrainUsed{} rows rather than
the full training period. Single-feature AUC moves with dataset size by more than seed noise, and the
cause is unexplained; every figure is therefore stated at its scale. The two failing gate metrics
(recall and false-negative rate at the flag threshold) may depend on training volume; that is a
candidate cause, not a finding.

\paragraph{Latency.} No latency figure here is a gate figure. The frontier is a median on a shared
development laptop, where the 99th percentile is not measurable; serving, end-to-end and load
latencies belong to later milestones.

\paragraph{Small counts.} Variant, country and channel results rest on tens to hundreds of fraud
rows; intervals are reported for each, and the Democratic Republic of the Congo's rest on fewer than
30 and give direction only.

\paragraph{Provisional results.} Every figure marked \textsuperscript{\dag} comes from a milestone not
yet tagged complete and is re-verified before publication. Two of the contribution figures already
moved between data draws (the redundancy threshold and the leave-one-country-out loss); the findings
did not.

\paragraph{Claims not made.} The specification's figures on regional fraud growth, payment volume,
analyst dismissal rates, the accuracy of models built elsewhere on East African data, annual model
degradation and ensemble variance reduction are unverified, replaced or refuted in the project's
claims register, and none is stated here as fact.
